\documentclass[11pt,a4paper]{article}
\input{../template/utm-lab}

\newcommand{\labnumber}{4}
\newcommand{\labtitle}{The xv6 Shell from the Outside}
\newcommand{\studentname}{Daniel Nicorici}
\newcommand{\studentgroup}{FAF-243}

\begin{document}

\utmtitlepage
\tableofcontents
\newpage

\plainsection{Purpose of the work}

The whole user interface of xv6 is one program, \code{user/sh.c}. The purpose of this lab is to
use it side by side with bash and find out, by experiment, what a shell must do and what bash
adds on top; to write a small program that prints the arguments it receives, which shows
exactly who expands \code{*} and who removes quotes; and to read the parts of \code{sh.c} that
explain what was observed. First, the \code{sleep} program from Lab 2 is checked against the
current xv6.

The work was done on my own laptop running Arch Linux: xv6 (\code{mit-pdos/xv6-riscv}, commit
\code{06aad25}) cross-compiled with \code{riscv64-linux-gnu-gcc} 16.2.1 and run on QEMU
11.1.1, and on the Linux side GNU bash 5.3.15 and gcc 16.1.1.

\section{Part 1. Fix sleep for the current xv6}

In August 2025 the xv6 authors renamed the system call \code{sleep()} to \code{pause()}. My
tree declares only \code{pause}, and my \code{sleep.c} from Lab 2 already calls \code{pause()},
so it builds without changes; \code{\$U/\_sleep\textbackslash} is on line 153 of the Makefile.

\begin{lstlisting}[caption={The grep result in user.h}]
$ cd ~/xv6-riscv
$ grep -n "pause\|sleep" user/user.h
25:int pause(int);
\end{lstlisting}

\begin{lstlisting}[style=ccode,caption={The updated sleep.c}]
#include "kernel/types.h"
#include "user/user.h"

int
main(int argc, char *argv[])
{
  if(argc != 2){
    fprintf(2, "usage: sleep <ticks>\n");
    exit(1);
  }
  pause(atoi(argv[1]));
  exit(0);
}
\end{lstlisting}

\needspace{17\baselineskip}
\begin{lstlisting}[caption={sleep running in xv6, with and without an argument}]
$ make qemu
qemu-system-riscv64 -machine virt -bios none -kernel kernel/kernel -m 128M -smp 3 -nographic -global virtio-mmio.force-legacy=false -drive file=fs.img,if=none,format=raw,id=x0 -device virtio-blk-device,drive=x0,bus=virtio-mmio-bus.0

xv6 kernel is booting

hart 2 starting
hart 1 starting
init: starting sh
$ sleep 20
$ sleep
usage: sleep <ticks>
\end{lstlisting}

\code{sleep 20} pauses for 20 ticks (about 2 seconds) before the prompt returns; without an
argument the program prints its usage message to standard error and exits with status 1.

\section{Part 2. bash versus the xv6 shell}

Every command was run twice: in bash, in an empty folder \code{\textasciitilde/shtest} containing
only \code{README}, and at the xv6 prompt. Table 2.1 summarises what each shell printed.

\begin{table}[H]
  \centering
  \caption{The same commands in bash and in the xv6 shell}
  \footnotesize
  \begin{tabular}{@{}>{\raggedright\arraybackslash}p{0.25\textwidth}>{\raggedright\arraybackslash}p{0.13\textwidth}>{\raggedright\arraybackslash}p{0.25\textwidth}>{\raggedright\arraybackslash}p{0.29\textwidth}@{}}
    \toprule
    \textbf{Command} & \textbf{Tests} & \textbf{bash} & \textbf{xv6} \\
    \midrule
    \code{echo hello > f ; cat < f} & redirection & \code{hello} & \code{hello} \\
    \code{ls | grep READ} & pipe & \code{README} & \code{README 2 2 2441} \\
    \code{(echo a; echo b) | wc} & grouping, list & \code{2 2 4} & \code{2 2 4} \\
    \code{echo one ; echo two} & sequence & \code{one}, \code{two} & \code{one}, \code{two} \\
    \code{echo bg \&} & background & \code{[1] 135185}, \code{bg}, later \code{[1]+ Done} & \code{bg}, no job number \\
    \code{ls *} & globbing & \code{README  f} & \code{ls: cannot open *} \\
    \code{echo \$HOME} & variables & \code{/home/Nico} & \code{\$HOME} \\
    \code{echo "a b"} & quotes & \code{a b} & \code{"a b"} \\
    \code{mkdir d ; cd d ; ls ..} & PATH & \code{README  d  f} & \code{exec cd failed}, then \code{ls} lists \code{/}; after a separate \code{cd d}: \code{exec ls failed} \\
    arrow up, Tab, Ctrl+R & history, completion & recalls, completes, searches & \code{exec \textasciicircum{}[[A failed}, \code{exec ech failed}, \code{exec \textasciicircum{}Rtwo failed} \\
    \bottomrule
  \end{tabular}
\end{table}

\begin{lstlisting}[caption={The commands run in bash}]
$ mkdir ~/shtest && cd ~/shtest && touch README
$ echo hello > f ; cat < f
hello
$ ls | grep READ
README
$ (echo a; echo b) | wc
      2       2       4
$ echo one ; echo two
one
two
$ echo bg &
[1] 135185
$ bg
ls *
README  f
[1]+  Done                       echo bg
$ echo $HOME
/home/Nico
$ echo "a b"
a b
$ mkdir d ; cd d ; ls ..
README  d  f
$ mkdir d ; cd d ; ls ..
d
$ echo  one
one
(reverse-i-search)`two': echo one ; echo two
one
two
\end{lstlisting}

In bash the background \code{echo} printed \code{bg} only after the next prompt had appeared,
so it shows up after \code{\$}, and the job's \code{Done} line came after the following
command. The last three commands test the keyboard: the arrow-up key recalled the previous line,
and running it again created \code{d/d} and listed it; typing \code{ech} and Tab completed the
word to \code{echo} plus a space; Ctrl+R followed by \code{two} found
\code{echo one ; echo two} in the history and ran it again.

\begin{lstlisting}[caption={The same commands run in xv6}]
$ echo hello > f ; cat < f
hello
$ ls | grep READ
README         2 2 2441
$ (echo a; echo b) | wc
2 2 4
$ echo one ; echo two
one
two
$ echo bg &
bg
$ ls *
ls: cannot open *
$ echo $HOME
$HOME
$ echo "a b"
"a b"
$ ^[[A
exec ^[[A failed
$ ech[Tab] one
exec ech failed
$ ^Rtwo
exec ^Rtwo failed
\end{lstlisting}

In xv6 the special keys are not interpreted: the arrow key arrives as the raw characters
\code{ESC [ A} (shown as \code{\textasciicircum{}[[A}), Tab is just a space-like character that ends the word
\code{ech}, and Ctrl+R is the raw character \code{\textasciicircum{}R}. The shell treats each of them as part
of a program name and fails to run it.

\begin{lstlisting}[caption={ls failing after cd d, and the working alternative}]
$ mkdir d ; cd d ; ls ..
exec cd failed
.              1 1 1024
..             1 1 1024
README         2 2 2441
[... 23 entries omitted ...]
f              2 26 6
d              1 27 32
$ cd d
$ ls ..
exec ls failed
$ /ls ..
.              1 1 1024
..             1 1 1024
README         2 2 2441
[... 23 entries omitted ...]
f              2 26 6
d              1 27 32
$ cd /
\end{lstlisting}

\textbf{Observations:}
\begin{enumerate}
  \item \textbf{Common features and bash-only features.} Both shells run programs with
  arguments, redirect input and output with \code{<} and \code{>}, connect programs with
  \code{|}, run sequences with \code{;}, group commands with \code{( )} and start background
  commands with \code{\&}. Only bash has globbing, variables, quoting, a PATH search, \code{cd}
  anywhere on the line, job numbers and \code{Done} messages, command history, Tab completion
  and Ctrl+R search.
  \item \textbf{Why \code{ls} stops working after \code{cd d}.} The xv6 shell passes the
  program name straight to \code{exec}, which looks it up relative to the current directory,
  and there is no PATH. All xv6 programs live in the root directory \code{/}, so after
  \code{cd d} the name \code{ls} no longer points to a file and the shell prints
  \code{exec ls failed}. The fix is to give the full path, \code{/ls ..} (or \code{../ls ..}).
  The first line also shows a second limitation: \code{cd} inside a \code{;} list is not
  recognised and fails with \code{exec cd failed} (explained in Part 4).
  \item \textbf{What I would miss most, and what is easiest to add.} In daily work I would miss
  history and Tab completion the most, because they save the most typing. The easiest feature
  to add is a PATH-like fallback: when \code{exec} fails and the name contains no \code{/}, the
  shell could retry with \code{/} in front of the name, which takes only a few lines in
  \code{runcmd}.
\end{enumerate}

\section{Part 3. Who expands *? A program that prints its argv}

\begin{lstlisting}[caption={args.c and the UPROGS line}]
$ cat user/args.c
#include "kernel/types.h"
#include "user/user.h"

int
main(int argc, char *argv[])
{
  printf("argc = %d\n", argc);
  for(int i = 0; i < argc; i++)
    printf("argv[%d] = [%s]\n", i, argv[i]);
  exit(0);
}
$ grep -n "_sleep\|_args" Makefile
153:	$U/_sleep\
154:	$U/_args\
\end{lstlisting}

\begin{lstlisting}[caption={args output in xv6}]
$ args one              two         three
argc = 4
argv[0] = [args]
argv[1] = [one]
argv[2] = [two]
argv[3] = [three]
$ args *.c "x y"
argc = 4
argv[0] = [args]
argv[1] = [*.c]
argv[2] = ["x]
argv[3] = [y"]
$ args $HOME
argc = 2
argv[0] = [args]
argv[1] = [$HOME]
$ args 'single quoted'
argc = 3
argv[0] = [args]
argv[1] = ['single]
argv[2] = [quoted']
\end{lstlisting}

\begin{lstlisting}[caption={args\_linux output in bash}]
$ cat args_linux.c
#include <stdio.h>

int
main(int argc, char *argv[])
{
  printf("argc = %d\n", argc);
  for(int i = 0; i < argc; i++)
    printf("argv[%d] = [%s]\n", i, argv[i]);
  return 0;
}
$ gcc -Wall -o args args_linux.c
$ ls *.c
args.c  args_linux.c  sleep.c
$ ./args one              two         three
argc = 4
argv[0] = [./args]
argv[1] = [one]
argv[2] = [two]
argv[3] = [three]
$ ./args *.c "x y"
argc = 5
argv[0] = [./args]
argv[1] = [args.c]
argv[2] = [args_linux.c]
argv[3] = [sleep.c]
argv[4] = [x y]
$ ./args $HOME
argc = 2
argv[0] = [./args]
argv[1] = [/home/Nico]
$ ./args 'single quoted'
argc = 2
argv[0] = [./args]
argv[1] = [single quoted]
\end{lstlisting}

\textbf{Observations:}
\begin{enumerate}
  \item \textbf{\code{args *.c "x y"} in both shells.} In xv6 the program received 4 arguments:
  \code{args}, \code{*.c}, \code{"x} and \code{y"}. In bash it received 5: \code{./args}, the
  three file names \code{args.c}, \code{args\_linux.c}, \code{sleep.c}, and \code{x y} as a
  single argument without quotes.
  \item \textbf{Who expands \code{*} and removes quotes.} The shell. The program is the same in
  both cases and only prints what it is given, and the kernel's \code{exec} only copies the
  argument strings it receives. Bash replaced \code{*.c} with the matching names and removed the
  quotes before calling \code{exec}; the xv6 shell does neither, so the program received the
  raw text. The same happened with \code{\$HOME} (expanded only by bash) and
  \code{\textquotesingle{}single quoted\textquotesingle{}} (one argument in bash, two in xv6).
  \item \textbf{The extra spaces in \code{one two three}.} They disappeared in both shells:
  both split the line at runs of whitespace, so spaces only separate arguments and are never
  part of them, and both programs received \code{argc = 4}.
\end{enumerate}

\section{Part 4. Read the shell}

\begin{lstlisting}[caption={The command kinds defined in sh.c}]
$ wc -l user/sh.c
499 user/sh.c
$ grep -n "#define" user/sh.c
8:#define EXEC  1
9:#define REDIR 2
10:#define PIPE  3
11:#define LIST  4
12:#define BACK  5
14:#define MAXARGS 10
$ grep -n "symbols\[\]\|whitespace\[\]" user/sh.c
268:char whitespace[] = " \t\r\n\v";
269:char symbols[] = "<|>&;()";
\end{lstlisting}

\begin{lstlisting}[caption={Where each kind is run}]
$ grep -n "case " user/sh.c
75:  case EXEC:
83:  case REDIR:
93:  case LIST:
101:  case PIPE:
125:  case BACK:
284:  case 0:
286:  case '|':
287:  case '(':
288:  case ')':
289:  case ';':
290:  case '&':
291:  case '<':
294:  case '>':
391:    case '<':
394:    case '>':
397:    case '+': // >>
469:  case EXEC:
475:  case REDIR:
481:  case PIPE:
487:  case LIST:
493:  case BACK:
\end{lstlisting}

Lines 75--125 are \code{runcmd()}, which executes each kind; lines 284--294 are the tokenizer
\code{gettoken()}, which recognises the special characters; lines 391--397 parse
redirections; and lines 469--493 are \code{nulterminate()}. Searching for \code{cd} finds only
the error message on line 170, because the check compares single characters:

\begin{lstlisting}
$ grep -n "cd" user/sh.c | head
170:        fprintf(2, "cannot cd %s\n", cmd + 3);
\end{lstlisting}

\begin{lstlisting}[style=ccode,firstnumber=159,caption={The code that handles cd}]
  // Read and run input commands.
  while (getcmd(buf, sizeof(buf)) >= 0) {
    char *cmd = buf;
    while (*cmd == ' ' || *cmd == '\t')
      cmd++;
    if (*cmd == '\n') // is a blank command
      continue;
    if (cmd[0] == 'c' && cmd[1] == 'd' && cmd[2] == ' ') {
      // Chdir must be called by the parent, not the child.
      cmd[strlen(cmd) - 1] = 0; // chop \n
      if (chdir(cmd + 3) < 0)
        fprintf(2, "cannot cd %s\n", cmd + 3);
    } else {
      if (fork1() == 0)
        runcmd(parsecmd(cmd));
      wait(0);
    }
  }
\end{lstlisting}

\begin{lstlisting}[caption={Searching sh.c for code that handles * or \$}]
$ grep -n "'\*'\|'\$'\|glob\|getenv" user/sh.c || echo "(no matches)"
(no matches)
\end{lstlisting}

\textbf{Observations:}
\begin{enumerate}
  \item \textbf{The five kinds of command.} \code{EXEC} is a plain command, produced by
  ordinary words with no special character (\code{ls ..}); \code{REDIR} is produced by
  \code{<}, \code{>} and \code{>>} (which the tokenizer turns into \code{+}); \code{PIPE} by
  \code{|}; \code{LIST} by \code{;}; and \code{BACK} by \code{\&}. Parentheses only group
  commands and are not a kind of their own.
  \item \textbf{Where \code{cd} is handled, and why not as a program.} In \code{main()}, lines
  166--171: if a line starts with \code{cd }, the shell itself calls \code{chdir} and does not
  fork. The current directory belongs to a process, so if a child process changed directory,
  the change would disappear when the child exits and the shell would stay where it was. Because
  the check only looks at the start of the line, \code{mkdir d ; cd d ; ls ..} goes to
  \code{runcmd()} instead, where \code{cd} is treated as a program and fails with
  \code{exec cd failed}, exactly as seen in Part 2.
  \item \textbf{Code for \code{*} or \code{\$}.} There is none. The only special characters
  the shell knows are \code{symbols[] = "<|>\&;()"} plus whitespace. This explains Part 2:
  \code{ls *} received a literal \code{*}, \code{echo \$HOME} printed \code{\$HOME}, and
  quotes stayed inside the arguments because they are ordinary characters to this shell.
\end{enumerate}

\plainsection{Conclusions}

Running the same commands in two shells separated what a shell must do from what bash adds for
comfort. The essential part is small: read a line, split it into words, recognise a handful of
special characters, and turn them into the five kinds of command that \code{sh.c} defines,
which it then runs with \code{fork}, \code{exec}, \code{pipe}, \code{dup} and \code{chdir}.
In 499 lines of C, xv6 already supports programs, redirection, pipes, sequences, grouping and
background jobs.

Everything else in bash, such as globbing, variables, quotes, PATH search, history and
completion, is convenience implemented inside the shell, not in the programs and not in the
kernel. The \code{args} program proved this directly: the same code printed \code{*.c} and the
quotes literally under xv6, and the expanded file names under bash. The experiment with
\code{cd} showed the one thing a shell cannot delegate to a child process, because the current
directory belongs to the process that changes it. Finally, the renaming of \code{sleep} to
\code{pause} was a reminder that the interface between user programs and the kernel is defined
by \code{user.h}, and code has to follow it when it changes.

\end{document}
